% Gemini theme
% https://github.com/anishathalye/gemini

\documentclass[final]{beamer}

% ====================
% Packages
% ====================

\usepackage[T1]{fontenc}
\usepackage{lmodern}
\usepackage[size=custom,width=118.9,height=84.1,scale=1.0]{beamerposter}
\usetheme{gemini}
\usecolortheme{gemini}
\usepackage{graphicx}
\usepackage{booktabs}
\usepackage{tikz}
\usepackage{pgfplots}
\pgfplotsset{compat=1.14}
\usepackage{anyfontsize}
\usepackage[inkscapelatex=false]{svg}

% ====================
% Lengths
% ====================

% If you have N columns, choose \sepwidth and \colwidth such that
% (N+1)*\sepwidth + N*\colwidth = \paperwidth
\newlength{\sepwidth}
\newlength{\colwidth}
\setlength{\sepwidth}{0.025\paperwidth}
\setlength{\colwidth}{0.3\paperwidth}

\newcommand{\separatorcolumn}{\begin{column}{\sepwidth}\end{column}}

% ====================
% Title
% ====================

\title{Model Compression API for a fast and precise inference in ExecuTorch}

\author{Daniil Lyakhov \inst{1} \and Aamir Nazir \inst{1}}

\institute[shortinst]{\inst{1} Intel Corporation}

% ====================
% Footer (optional)
% ====================

\footercontent{
  \href{https://www.example.com}{https://www.example.com} \hfill
  ABC Conference 2025, New York --- XYZ-1234 \hfill
  \href{mailto:alyssa.p.hacker@example.com}{alyssa.p.hacker@example.com}}
% (can be left out to remove footer)

% ====================
% Logo (optional)
% ====================

% use this to include logos on the left and/or right side of the header:
% \logoright{\includegraphics[height=7cm]{logo1.pdf}}
% \logoleft{\includegraphics[height=7cm]{logo2.pdf}}

% ====================
% Body
% ====================

\begin{document}

\begin{frame}[t]
  \begin{columns}[t]
    \separatorcolumn

    \begin{column}{\colwidth}

      \begin{block}{Abstract}

        We present a \textbf{Model Compression API} for ExecuTorch, powered by the \textbf{Neural Network Compression Framework (NNCF)}.
        It extends the standard \texttt{torch.ao} quantization pipeline with data-aware algorithms --- \texttt{quantize\_pt2e} for post-training quantization and \texttt{compress\_pt2e} for weight compression.
        \par

        We demonstrate that the model compression with the \textbf{NNCF API} yields significant speedups on Llama~3.2, YOLOv12, and Stable Diffusion while preserving target metrics --- for the \textbf{OpenVINO ExecuTorch backend}.
        \par

        The NNCF compression API operates on the quantization annotations provided by a given \texttt{torch.ao} Quantizer --- rather than on backend-specific internals --- so any vendor backend can leverage its data-aware algorithms without altering the compressed model architecture.

      \end{block}

      \begin{block}{torch.ao pipeline vs NNCF API}

        NNCF's \texttt{quantize\_pt2e} and \texttt{compress\_pt2e} extend the standart torch.ao quantization pipeline with data-aware algorithms while preserving compatibility with the \texttt{torch.ao} Quantizer.

        \begin{figure}
          \centering
          \includesvg[width=0.95\colwidth]{figures/quantization_flow}
          \caption{Standard torch.ao PTQ flow (left) vs.\ NNCF \texttt{quantize\_pt2e} / \texttt{compress\_pt2e} (right).}
        \end{figure}

      \end{block}

      \begin{block}{Post-Training Quantization --- \texttt{quantize\_pt2e}}

        Can be applied on top of \emph{any} \texttt{torch.ao} Quantizer to improve the accuracy of the quantized models. The API implements the standard \texttt{prepare\_pt2e}/\texttt{convert\_pt2e} flow and augments it with data-aware algorithms:

        \begin{enumerate}
          \item \textbf{SmoothQuant} --- Reduces activation quantization error by inserting smoothing scales before weighted layers~\cite{smoothquant}.
          \item \textbf{BiasCorrection} --- Corrects bias shifts introduced by quantization~\cite{bias_correction}.
        \end{enumerate}

        This means existing quantization pipelines can benefit from NNCF's accuracy-boosting algorithms without changing the quantizer --- simply replace the standard \texttt{torch.ao} quantization call with \texttt{quantize\_pt2e} and supply a calibration dataset.


      \end{block}



    \end{column}

    \begin{column}{\colwidth}

      \begin{block}{Weight Compression --- \texttt{compress\_pt2e}}

        \heading{}

        Applies weight compression to a \texttt{torch.fx.GraphModule}, targeting LLM deployment. Key algorithms:

        \begin{itemize}
          \item \textbf{AWQ}  --- Activation-aware Weight Quantization that finds optimal per-channel scales to minimize quantization error based on activation distributions~\cite{awq}.
          \item \textbf{Scale Estimation}  --- Estimates optimal quantization scales for 4-bit layers to further reduce perplexity of compressed LLMs.
          \item \textbf{GPTQ}  --- Applies the GPTQ algorithm for accurate layer-wise weight quantization using second-order information~\cite{gptq}.
          \item \textbf{LoRA Correction}  --- Applies Low-Rank Adaptation corrections to compensate for weight compression errors~\cite{lora}.
        \end{itemize}

        Currently, \texttt{compress\_pt2e} supports the \texttt{OpenVINOQuantizer} and the OpenVINO backend. Support for third-party quantizers can be added by implementing a lightweight NNCF quantizer adapter.

      \end{block}

      \begin{block}{OpenVINO specific quantizer - OpenVINOQuantizer}

        The \texttt{OpenVINOQuantizer} from NNCF is a \texttt{torch.ao} Quantizer tailored for the OpenVINO backend. It unlocks the full potential of low-precision OpenVINO kernels by placing quantizers designed specifically for OpenVINO operator implementations.

        \heading{Configuration}

        The quantizer is configured via a single \textbf{\texttt{mode}} parameter of type \texttt{QuantizationMode}:

        \begin{itemize}
          \item \textbf{\texttt{INT8\_SYM}} (default) --- INT8 symmetric quantization for both activations and weights.
          \item \textbf{\texttt{INT8\_MIXED}} --- INT8 asymmetric quantization for activations, symmetric for weights.
          \item \textbf{\texttt{INT8\_TRANSFORMER}} --- Optimized INT8 quantization for Transformer-based architectures (BERT, Llama, etc.).
          \item \textbf{\texttt{INT4WO\_SYM / INT4WO\_ASYM}} --- INT4 symmetric/asymmetric weights-only compression.
          \item \textbf{\texttt{INT8WO\_SYM / INT8WO\_ASYM}} --- INT8 symmetric/asymmetric weights-only compression.
        \end{itemize}

        Additionally, \texttt{set\_ignored\_scope} allows excluding specific layers from quantization by node names, regex patterns, operation types, or subgraphs --- useful for preserving accuracy-sensitive parts of a model.

      \end{block}

      \begin{block}{Results}

        \heading{LLM: Llama~3.2 Weight Compression}

        \begin{table}
          \centering
          \begin{tabular}{l r r}
            \toprule
            \textbf{Model / Mode}            & \textbf{Perplexity $\downarrow$} & \textbf{Model Size} \\
            \midrule
            Llama~3.2 1B (FP16)              & ---                              & ---                 \\
            Llama~3.2 1B (INT8)              & ---                              & ---                 \\
            Llama~3.2 1B (INT4 + Scale Est.) & ---                              & ---                 \\
            \midrule
            Llama~3.2 3B (FP16)              & ---                              & ---                 \\
            Llama~3.2 3B (INT8)              & ---                              & ---                 \\
            Llama~3.2 3B (INT4 + Scale Est.) & ---                              & ---                 \\
            \bottomrule
          \end{tabular}
          \caption{Llama~3.2 perplexity and model size under different compression modes.}
        \end{table}

        \heading{Detection: YOLOv12 INT8 Quantization}

        INT8 post-training quantization applied via \texttt{quantize\_pt2e} with SmoothQuant and BiasCorrection~\cite{yolo_et}.

        \begin{table}
          \centering
          \begin{tabular}{l r r r}
            \toprule
            \textbf{Model / Mode} & \textbf{mAP@50 $\uparrow$} & \textbf{mAP@50:95 $\uparrow$} & \textbf{Model Size} \\
            \midrule
            YOLOv12 (FP32)        & ---                        & ---                           & ---                 \\
            YOLOv12 (INT8)        & ---                        & ---                           & ---                 \\
            \bottomrule
          \end{tabular}
          \caption{YOLOv12 detection accuracy: FP32 baseline vs.\ INT8 quantized.}
        \end{table}
      \end{block}
    \end{column}

    \separatorcolumn

    \begin{column}{\colwidth}



      \heading{Image Generation: Stable Diffusion}

      Mixed-precision pipeline: INT8 quantization for U-Net via \texttt{quantize\_pt2e}, weight compression for text encoder and VAE decoder via \texttt{compress\_pt2e}~\cite{sd_et}.

      \begin{figure}
        \centering
        \begin{tabular}{ccc}
          % Pair 1
          \includegraphics[width=0.28\colwidth]{example-image} &
          \includegraphics[width=0.28\colwidth]{example-image} &
          \includegraphics[width=0.28\colwidth]{example-image}                                                           \\
          % Pair 2
          \includegraphics[width=0.28\colwidth]{example-image} &
          \includegraphics[width=0.28\colwidth]{example-image} &
          \includegraphics[width=0.28\colwidth]{example-image}                                                           \\
          {\small FP32}                                        & {\small FP32}              & {\small FP32}              \\[2pt]
          \includegraphics[width=0.28\colwidth]{example-image} &
          \includegraphics[width=0.28\colwidth]{example-image} &
          \includegraphics[width=0.28\colwidth]{example-image}                                                           \\
          {\small INT8 / Compressed}                           & {\small INT8 / Compressed} & {\small INT8 / Compressed} \\
        \end{tabular}
        % Replace 'example-image' with actual image paths, e.g.:
        % figures/sd_prompt1_fp32.png, figures/sd_prompt1_int8.png, etc.
        \caption{Stable Diffusion generation quality: FP32 (top) vs.\ INT8/compressed (bottom) for three prompts.}
      \end{figure}


      \begin{block}{References}

        \nocite{*}
        \footnotesize{\bibliographystyle{plain}\bibliography{poster}}

      \end{block}

    \end{column}

    \separatorcolumn
  \end{columns}
\end{frame}

\end{document}
